%% Methodology comparison of independent models of XUV-driven atmospheric
%% escape: AIOLOS, Taylor et al. (2025, 2026), Xing et al. (2023) and
%% Frelikh & Murray-Clay (2026), and what EXHALE has taken from each.  The
%% EXHALE column and the whole of the adoption assessment were re-checked
%% against src/ rather than against the documentation; the four external
%% columns are what those papers say, except where a number is marked as
%% measured here.  Source PDFs under ../../references/.
%% Uses the local my_memo document class (docs/my_memo.cls), which auto-loads
%% graphicx, amsmath, amssymb, hyperref, xcolor, kotex, fancyhdr, sets its own
%% A4 geometry / boxed title / header-footer, and defaults \date to the
%% modification date of this file.  booktabs, array, longtable and enumitem
%% are not provided by the class and are loaded here.
%% Author: Kwang-Il Seon
\documentclass[english,a4paper]{my_memo}
\synctex=1
\usepackage{booktabs}
\usepackage{array}
\usepackage{longtable}
\usepackage{enumitem}
\emergencystretch=3em

%% Ragged-right fixed-width column: the comparison table's columns are narrow
%% enough that justification opens large interword gaps.
\newcolumntype{L}[1]{>{\raggedright\arraybackslash}p{#1}}
%% Compact table body: the class sets \linespread{1.5} for running text, which
%% makes multi-line p-columns unnecessarily tall.
\newcommand{\tabsetup}{\footnotesize\linespread{1.05}\selectfont
  \renewcommand{\arraystretch}{1.15}}

\newcommand{\Rp}{R_{\rm p}}
\newcommand{\Mdot}{\dot{M}}
\newcommand{\EXHALE}{\textsc{Exhale}}
\newcommand{\AIOLOS}{\textsc{Aiolos}}
\newcommand{\Lya}{Ly$\alpha$}
\newcommand{\Ha}{H$\alpha$}
\newcommand{\Kzz}{K_{zz}}
%% Mode-safe: these appear both in running text and inside math expressions.
\newcommand{\Hth}{\ensuremath{\mathrm{H}_3^+}}
\newcommand{\Htwo}{\ensuremath{\mathrm{H}_2}}

\title{Methodology comparison: independent models of\\
XUV-driven atmospheric escape, and what \EXHALE\ took from each}
\author{Kwang-Il Seon}

\begin{document}
\maketitle

\section{Purpose, and the verdict first}

This memo compares the methodology of four independent modelling efforts for
XUV-driven atmospheric escape from close-in exoplanets: the general-purpose
1-D radiation-hydrodynamics code \AIOLOS\ (Schulik \& Booth 2023, MNRAS 523,
286); the dedicated thermosphere--ionosphere escape model of Taylor et al.\
(2025, 2026), built on the Koskinen et al.\ code; the multi-fluid PLUTO-based
model of Xing et al.\ (2023); and the photochemical--hydrodynamic model of
Frelikh \& Murray-Clay (2026), and then records, item by item, what
\EXHALE\ has taken from each and what it has not.

The memo was first written as a forward-looking recommendation list. It is now
a record of outcomes, so every candidate carries one of three verdicts:
\textbf{adopted} (and in which module), \textbf{adopted then replaced} (and
why the first choice did not survive), or \textbf{not adopted} (and whether the
reason still holds). Of the fifteen items assessed in the adoption assessment
below, six are adopted, two were adopted and then replaced, and seven are not
adopted. The two largest gaps the original assessment identified (no
radius-dependent He/H separation, and an ad hoc lower boundary) are both
closed, and closed further than the recommendation asked for.

Frelikh \& Murray-Clay is the newest of the four and the closest peer to
\EXHALE's molecular layer, so it is the source of four of the seven
not-adopted items and of one finding that belongs to neither code alone: the
thermal-dissociation rate pair that both use is not thermodynamically
consistent below about 3000~K (Section~\ref{sec:detbal}). Three further points
where the two codes diverge, and where the judgement went against \EXHALE, are
the photoelectron partition in a molecular gas (Section~\ref{sec:secion}), the
absence of a composition boundary condition at the base
(Section~\ref{sec:lowerbc}), and the sharpness of the modeled
\Htwo$\rightarrow$H transition (Section~\ref{sec:mollayer}).

Every statement about \EXHALE\ below was checked against \texttt{src/} rather
than against \EXHALE's own documentation; the file and routine names given are
the ones in the tree. The \AIOLOS, Taylor, Xing and Frelikh entries are what
those papers say. Numbers measured here from \EXHALE\ output, rather than
quoted from a paper, are marked as such at the point of use.

\section{The codes considered}

\begin{itemize}[nosep]
  \item \textbf{\AIOLOS}: the general-purpose 1-D multi-species
    radiation-hydrodynamics code of Schulik \& Booth (2023, MNRAS 523, 286;
    ``SB23''), distributed at \url{github.com/Schulik/aiolos}.
  \item \textbf{Taylor et al.\ (2025)}, \emph{ApJ} 989, 68: ``A Multispecies
    Atmospheric Escape Model with Excited Hydrogen and Helium: Application to
    HD209458b,'' built on the Koskinen et al.\ (2013a,b; 2022)
    thermosphere--ionosphere code.
  \item \textbf{Taylor et al.\ (2026)}, \emph{ApJ} 999, 214: ``Helium Escape
    in Context: Comparative Signatures of Four Close-in Exoplanets,'' a direct
    application and extension of Taylor (2025) to HD~209458\,b,
    HD~189733\,b, HD~149026\,b and GJ~1214\,b.
  \item \textbf{Xing et al.\ (2023)}, \emph{ApJ} 953, 166: ``The Mass
    Fractionation of Helium in the Escaping Atmosphere of HD~209458b,'' a
    \emph{multi-fluid} model built inside PLUTO.
  \item \textbf{Frelikh \& Murray-Clay (2026)}, \emph{ApJ} 996, 96:
    ``Efficiency of Hydrodynamic Atmospheric Escape in Hot Jupiters and
    Super-Earths,'' a 1-D photochemical--hydrodynamic model of H--He escape
    with \Hth\ and \Lya\ cooling, heat conduction, tidal gravity and
    photoelectron secondary ionization.
\end{itemize}
``Taylor'' denotes the shared Koskinen-based framework used in both Taylor
papers unless a specific paper is named. ``Frelikh'' denotes Frelikh \&
Murray-Clay (2026).

\section{One-paragraph characterization of each code}

\paragraph{\AIOLOS.} A \emph{general-purpose} radiation-hydrodynamics code, not
a dedicated escape code. It solves the time-dependent Euler equations (mass,
momentum, energy) for an arbitrary number of species (each a
\textbf{separate fluid} with its own density, momentum and internal energy)
coupled by an inter-species \textbf{friction/drag} solver. It uses
finite-volume \textbf{HLLC} fluxes with \textbf{PLM} reconstruction
(monotonized-central limiter) and second-order
\textbf{strong-stability-preserving} Runge--Kutta time stepping, in
Cartesian, cylindrical or \textbf{spherical} geometry. Radiation is treated
with \textbf{flux-limited diffusion (FLD)} over arbitrary in/out spectral
bands, solved implicitly through a block-tridiagonal system. Escape is one of
many problems it can run.

Two things about \AIOLOS\ have to be attributed to the distributed code rather
than to SB23, because the paper does not carry them. The paper implements a
C2Ray-style two-species photoionization scheme and calls a fuller chemical
solver future work; the code has since acquired a reaction-network class
(\texttt{c\_reaction}) with selectable photochemistry levels, the Black (1981)
heating and cooling rates with free-free and H$_3^+$ cooling
(\texttt{photochem.cpp}, \texttt{chemistry.cpp}), grid spacing that can be
linear, logarithmic or bi-logarithmic (\texttt{PARI\_GRID\_TYPE}), and a
tidal-field option (\texttt{USE\_TIDES}). Statements below about \AIOLOS's
chemistry, cooling, grid or tides refer to the code.

\paragraph{Taylor (2025/2026).} A \emph{dedicated} 1-D spherical
thermosphere--ionosphere and hydrodynamic-escape model (Koskinen et al.\
2013/2022), integrated to steady state, with a \textbf{single bulk momentum
equation} but \textbf{multi-species molecular and eddy diffusion} that
redistributes species, so He and H can separate diffusively. Its distinguishing
strengths are in the completeness of the \emph{microphysics}: a full
He($2^3$S) and H($n{=}2$) non-LTE network solved by the KPP kinetic
preprocessor, a \textbf{\Lya\ Monte Carlo} radiative-transfer model (Huang et
al.\ 2017) iterated with the hydro solution for \Ha, a high-resolution
He($2^3$S) photoionization cross section (B-spline K-matrix),
temperature-dependent Penning ionization, and self-consistent coupling to a
photochemical lower and middle atmosphere at the $\mu$bar level.

\paragraph{Xing (2023).} A \emph{dedicated} 1-D spherical escape model whose
distinguishing feature is that it is genuinely \textbf{multi-fluid}: H, H$^+$,
He, He$^+$ and e$^-$ each carry their \textbf{own velocity}, coupled by
resonant and non-resonant collisional drag. It is the only one of the three
that captures \textbf{He/H mass fractionation} (heavy He under-dragged by
escaping H) in the dynamics itself rather than through a diffusion
coefficient. It is built on PLUTO (HLL Riemann solver, third-order TVD RK) with
a Riemann split that carries the electron/electric-field pressure term. The
He($2^3$S) 10830~\AA\ line is post-processed, not solved in the hydro loop.

\paragraph{Frelikh (2026).} A \emph{dedicated} 1-D spherical escape model
built around a molecular base, and of the four the closest peer to \EXHALE's
lower atmosphere. It is \textbf{single-fluid} in velocity (all species
travel with the bulk $u$, and the diffusion the code contains is turned off
for the published H--He runs), but it carries a mass-conservation equation
for each of the eight H and He species with a chemical source term. The
hydrodynamics are unusual for this field: no Riemann problem at all, but the
\textbf{constrained interpolation profile} (CIP) advection scheme of Yabe \&
Aoki (1991) on a \textbf{staggered mesh}, operator-split against explicitly
differenced source terms, with a Landshoff-plus-von~Neumann artificial
viscosity for stability. The 22-reaction H--He network (their Table~1,
compiled from Yelle 2004 and Garc\'ia-Mu\~noz 2007) is \textbf{integrated in
time} as a stiff ODE system by the semi-implicit Euler stepper
\texttt{StepperSie} of Press et al.\ (2007) with an \textbf{analytic
Jacobian}. Radiative cooling is \Lya\ (Black 1981) and \Hth\ (Miller et al.\
2013); heat conduction follows Garc\'ia-Mu\~noz (2007). Photoelectron
secondary ionization uses Dalgarno et al.\ (1999), including the \Htwo\
vibrational, triplet-dissociation and Lyman--Werner channels. The published
result is a three-layer atmosphere: an \Hth-cooled \Htwo\ layer at the base, a
\Lya-cooled neutral-H layer above it, and a PdV-cooled ionized outflow. There
is no He metastable, no He$\leftrightarrow$H charge exchange and no metal in
the network.

\paragraph{\EXHALE.} A 1-D spherical \textbf{single-fluid} escape code of ATES
heritage. The hydrodynamics are finite-volume, second-order, with PLM or WENO3
reconstruction (usable as a two-stage PLM$\rightarrow$WENO3 sequence) and a
selectable LLF/\textbf{HLLC}/Roe numerical flux
(\texttt{flux/Num\_Fluxes.f90}); reconstructed face states are dropped to the
piecewise-constant limit wherever they lose positivity
(\texttt{positivity\_limited\_faces}, \texttt{states/Reconstruction.f90}), and
a gated fourth-difference dissipation restores the damping of the contact mode
that HLLC loses as the flow stalls (\texttt{flux/low\_mach\_dissipation.f90},
key \texttt{Low-Mach damping}, off by default). The steady wind is reached either by explicit marching on the mass-flux
($\rho v r^2$) spread or by a direct \textbf{Jacobian-free Newton--Krylov}
solve: right-preconditioned GMRES($m$) with pseudo-transient continuation
(\texttt{time\_step/steady\_newton.f90}, key \texttt{Solver: Newton}).
Ionization of H, He, ten metal elements in 27 ion stages, the He($2^3$S)
metastable and, when molecular chemistry is on, H$_2$/H$_2^+$/H$_3^+$/HeH$^+$
is solved as \emph{one} coupled nonlinear algebraic system (analytic-Jacobian
Newton, MINPACK \texttt{hybrd1} fallback; the \texttt{System\_HeH\_*} family).
The photoelectron heating fraction is not a free parameter: it follows the
Shull \& van Steenberg (1985) branching evaluated on each cell's own ionized
fraction and neutral composition. Helium is transported as the second
component of a binary mixture with stage-resolved friction
(\texttt{functions/binary\_element\_diffusion.f90}), the Navier--Stokes viscous
force and heat conduction are available
(\texttt{time\_step/viscous\_conduction.f90}), and the base can be handed over
from a photochemical lower atmosphere as scalars or as a full profile.
Diagnostics are a non-LTE He($2^3$S) network, a non-LTE H($n{=}2$) population
(Christie et al.\ 2013), \Lya\ transfer through a Neufeld escape-probability
closure or an imported field, and the transmission post-processor
\texttt{EXHALE\_transit.py}. Core size: 75 source files, 16\,568 SLOC excluding
blank and comment lines (\texttt{src/utils/codesize.py}, measured 2026-08-30),
of which 6\,404 are in modules that do not exist in upstream ATES.

\section{Side-by-side methodology table}
\label{sec:table}

A sixth column does not fit across one A4 text block, so the comparison is
split by axis: Table~\ref{tab:methods-num} carries the numerical method and
Table~\ref{tab:methods-phys} the physics and the diagnostics.

{\tabsetup
\begin{longtable}{@{}L{0.12\textwidth} L{0.150\textwidth} L{0.150\textwidth}
                    L{0.150\textwidth} L{0.150\textwidth} L{0.150\textwidth}@{}}
\caption{Numerical method. The four external columns are as described in their
papers; the \EXHALE\ entries were verified in
\texttt{EXHALE\_v1.00/src/}.}\label{tab:methods-num}\\
\toprule
\textbf{Axis} & \textbf{\AIOLOS\ (SB23)} & \textbf{Taylor 2025/2026} &
\textbf{Xing 2023} & \textbf{Frelikh 2026} & \textbf{\EXHALE} \\
\midrule
\endfirsthead
\multicolumn{6}{@{}l}{\emph{Table \ref{tab:methods-num}, continued}}\\
\toprule
\textbf{Axis} & \textbf{\AIOLOS} & \textbf{Taylor} & \textbf{Xing} &
\textbf{Frelikh} & \textbf{\EXHALE} \\
\midrule
\endhead
\bottomrule
\endfoot
Purpose & General multi-species RHD & Dedicated escape $+$ line diagnostics &
Dedicated escape, He/H fractionation & Dedicated escape; radiative-cooling
efficiency & Dedicated escape $+$ line diagnostics \\
Dimensions & 1-D (sph./cyl./cart.) & 1-D spherical & 1-D spherical &
1-D spherical, substellar ray & 1-D spherical \\
Route to steady state & Time-dependent $\to$ steady & Time-dependent $\to$
steady & Time-dependent $\to$ steady & Time-dependent $\to$ steady &
RK marching on the $\rho v r^2$ spread, optionally finished by a direct JFNK
solve \\
Fluid model & \textbf{Multi-fluid} $+$ friction & \textbf{Single bulk} $+$
multi-species diffusion & \textbf{Multi-fluid} & \textbf{Single fluid}; one
mass equation per species, all at the bulk $u$ & \textbf{Single fluid} $+$
binary element diffusion \\
Advection scheme & PLM $+$ MC limiter & Flux-conservative van Leer advection
(Koskinen et al.\ 2013a); not restated in the Taylor papers & PLM (PLUTO) &
\textbf{CIP} cubic interpolation carrying $X$ and $\partial_r X$, on a
staggered mesh & PLM, WENO3, or two-stage PLM$\to$WENO3, with a face
positivity guard \\
Riemann solver & HLLC & - & HLL ($+$ e$^-$-pressure split) & \textbf{None};
Landshoff $+$ von~Neumann artificial viscosity & LLF, HLLC or Roe
(\texttt{Numerical flux}) \\
Time integrator & Second-order SSP RK & Operator split, semi-implicit
Crank--Nicholson, forward to steady & Third-order TVD RK & Operator split: CIP
advection, then explicitly differenced source terms & RK marching;
GMRES($m$) JFNK with PTC for the steady solve \\
Grid & Linear, log or bi-log (code) & Radial & 1024 points, stretched,
1--10\,$\Rp$ & Nonuniform, staggered, 1--8.8\,$\Rp$, ghost cells &
Uniform, stretched or mixed; cell count set at runtime \\
\end{longtable}}

{\tabsetup
\begin{longtable}{@{}L{0.12\textwidth} L{0.150\textwidth} L{0.150\textwidth}
                    L{0.150\textwidth} L{0.150\textwidth} L{0.150\textwidth}@{}}
\caption{Physics and diagnostics. Conventions as in
Table~\ref{tab:methods-num}. \EXHALE\ options marked \emph{opt-in} default to
off.}\label{tab:methods-phys}\\
\toprule
\textbf{Axis} & \textbf{\AIOLOS\ (SB23)} & \textbf{Taylor 2025/2026} &
\textbf{Xing 2023} & \textbf{Frelikh 2026} & \textbf{\EXHALE} \\
\midrule
\endfirsthead
\multicolumn{6}{@{}l}{\emph{Table \ref{tab:methods-phys}, continued}}\\
\toprule
\textbf{Axis} & \textbf{\AIOLOS} & \textbf{Taylor} & \textbf{Xing} &
\textbf{Frelikh} & \textbf{\EXHALE} \\
\midrule
\endhead
\bottomrule
\endfoot
Radiation & \textbf{FLD}, multi-band, implicit & XUV attenuation $+$ \Lya\
Monte Carlo & Radial XUV attenuation, 53 bins & Radial XUV attenuation, 37
bins (Richards et al.\ 1994), 12--248\,eV & Radial XUV attenuation on a
multi-bin energy grid (power law or tabulated SED) $+$ \Lya\ transfer \\
Photoelectron heating & From the C2Ray energy split & Tuned efficiency,
20--40\,\% & Tuned frequency-averaged $\eta$, 0.1--0.5 & \textbf{Dalgarno et
al.\ (1999)} $W_0(1{+}Cx^\alpha)$, with the \Htwo\ vibrational, triplet and
Lyman--Werner returns added back & Shull \& van Steenberg (1985) branching
renormalized on the cell's own neutral composition; \textbf{no efficiency
parameter} \\
Chemistry solve & C2Ray two-species (paper); reaction network (code) &
KPP preprocessor over a stiff ODE network & Semi-implicit source terms &
\textbf{Stiff ODE integrated in time}: semi-implicit Euler
(\texttt{StepperSie}), \textbf{analytic} Jacobian, LU & \textbf{One coupled
nonlinear algebraic system} per cell; analytic-Jacobian Newton, or MINPACK
\texttt{hybrd} under a continuation \\
He/H separation & Via friction/drag & Molecular $+$ eddy diffusion of the
trace ratio $n_{\rm He}/n_{\rm H}$ & Via multi-fluid dynamics & Implemented
but turned off for the published H--He runs & Binary diffusion of the mass
fraction $X{=}\rho_{\rm He}/\rho$, friction resolved by ionization stage;
\emph{opt-in} \\
Metals & Through the code's chemistry & Optional (solar) & None & \textbf{None}
& \textbf{10 elements, 27 ion stages} in the coupled system, with CHIANTI
closed-form line cooling \\
Molecules & Through the code's chemistry & \Htwo, H$_2^+$, \Hth, HeH$^+$
(2026, for HD~189733\,b and GJ~1214\,b only) & None & Same four species; 22
reactions (their Table~1) & Same four species inside the coupled system, with
Lyman--Werner photodissociation; \emph{opt-in} \\
He$\leftrightarrow$H charge exchange & Through the code's chemistry & Yes &
Yes & \textbf{Absent from the network} & Huang et al.\ (2023) Table~4, rows
B1/B2: a non-radiative and a radiative channel, separately \\
Cooling & Recombination $+$ \Lya\ (paper); Black (1981) line/free-free $+$
\Hth\ (code) & Recombination $+$ H\,\textsc{i} lines $+$ \Hth\ &
\textbf{\Lya\ only} & \Lya\ $+$ \Hth\ only; recombination cooling checked and
left off & \Lya, recombination, free-free, CHIANTI metal lines with
fine-structure line trapping, \Hth\ infrared \\
\Hth\ cooling form & Black (1981) (code) & Miller et al.\ (2013) & - &
Miller et al.\ (2013), LTE, emission into vacuum; zero above 5000\,K &
Miller et al.\ (2013) with the non-LTE departure factor; net exchange with a
diluted base field, \emph{opt-in} \\
He($2^3$S) 10830 & Not a focus & In-loop non-LTE & Post-processed
(Yan et al.\ 2022) & \textbf{Not in the network} & In-loop non-LTE (default
on) $+$ \texttt{EXHALE\_transit.py} \\
H($n{=}2$)/\Ha & Not a focus & Non-LTE $+$ \Lya\ Monte Carlo & Post-processed &
Not computed & Non-LTE (Christie et al.\ 2013) $+$ \Lya\ escape probability or
an imported field \\
Viscosity / conduction & - & Molecular $+$ eddy diffusion, conduction &
Collisional drag & Heat conduction (Garc\'ia-Mu\~noz 2007); no viscosity &
Navier--Stokes viscosity and heat conduction (Crank--Nicolson); \emph{opt-in}
\\
Tidal/Roche & \texttt{USE\_TIDES} (code) & Tested (2025); off (2026) & Stellar
tidal term & Tidal term (their Appendix~C) & Roche potential, spherical or
Roche-lobe domain \\
Lower boundary & Configurable & \textbf{$\mu$bar, coupled to photochemistry} &
1\,$\Rp$, fixed $n$, $T = 1500$\,K & $\approx$1\,$\mu$bar: \textbf{species
densities and $p$ imposed}, $u$ extrapolated; Newtonian relaxation to $T_0$
below $\tau{=}3$ & $\mu$bar base: fixed $T$,$n$; or an analytic column; or a
scalar \texttt{base.inp}; or a full photochemical profile with $\Kzz(p)$ \\
Transmission & - & Ray-traced Voigt & Post-processed 10830 & \Hth\ emission
spectrum (ExoCross); no transit spectrum &
\texttt{EXHALE\_transit.py}: 10830, \Lya, \Ha, H$\beta$, metal doublets,
optional triaxial Roche geometry \\
\end{longtable}}

\section{What they have in common}
\label{sec:common}

\begin{enumerate}[nosep]
  \item \textbf{1-D spherical XUV-driven escape}, with photoionization heating
    balanced by \Lya/recombination cooling and adiabatic expansion, is the
    shared physical backbone.
  \item \textbf{A steady transonic wind is the target.} The four external
    codes evolve the time-dependent equations forward; \EXHALE\ marches in the
    same way but can also solve the steady equations directly. The physical end
    state is the same.
  \item \textbf{Much of the same core H/He atomic physics.} Voronov collisional
    ionization, radiative and dielectronic recombination, and
    H$\leftrightarrow$He charge exchange appear in Xing, Taylor and \EXHALE\
    with near-identical rate coefficients, differing mainly in the source table
    (Storey \& Hummer, Badnell, Verner). \EXHALE's charge-exchange network is
    now Huang et al.\ (2023) Table 4 (\texttt{radiation/charge\_exchange.f90}).
    Frelikh is the exception on the last of these: its Table~1 has no
    He$^+ + $H and no He$ + $H$^+$ reaction of any kind, and neither does its
    text (Section~\ref{sec:network}).
  \item \textbf{He($2^3$S) as the key observable.} Taylor, Xing and \EXHALE\
    all target the 10830~\AA\ line through a non-LTE $2^3$S population
    (recombination, collisions, radiative decay, Penning ionization,
    photoionization). \AIOLOS\ and Frelikh do not; the He metastable is not a
    species in either.
  \item \textbf{Finite-volume, slope-limited, shock-capturing hydrodynamics}
    underlies \AIOLOS, Xing (PLUTO) and \EXHALE. Two of the five are outliers,
    and in the same direction: the Koskinen solver beneath Taylor is
    flux-conservative van Leer advection operator-split against a semi-implicit
    Crank--Nicholson step (Koskinen et al.\ 2013a), and Frelikh is CIP
    advection operator-split against explicitly differenced source terms. Both
    are staggered, both use an added dissipation in place of an upwind flux,
    and \textbf{neither poses a Riemann problem}. Neither Taylor paper restates
    its numerical scheme; Frelikh gives its scheme in full in its Appendix~A.
\end{enumerate}

One item that used to be on this list no longer belongs. It read: none of the
four derives the photoelectron heating efficiency from first principles, and
all treat it as the primary free knob. That is no longer true of \EXHALE. Its
heating rate is the integral of the excess photon energy over the energy grid
weighted by the Shull \& van Steenberg (1985) heat fraction, and that fraction
is a function of the local ionized fraction with the ionization branching
renormalized onto the cell's own neutral H and He
(\texttt{svs85\_secondary\_branching}, \texttt{radiation/util\_ion\_eq.f90});
there is no heating-efficiency key in \texttt{input.inp} at all. Taylor and
Xing still tune theirs, and \AIOLOS\ takes its split from C2Ray.

\section{Where they genuinely differ}

\subsection{Fluid model}
\label{sec:fluid}
\begin{itemize}[nosep]
  \item \textbf{Xing is multi-fluid}, each species carrying its own velocity,
    which is what reproduces \textbf{He/H mass fractionation}: because He is
    four times heavier it is under-dragged by the H wind, so the
    \emph{escaping} He/H is about half the bulk abundance, a weak 10830 line
    without a subsolar He abundance.
  \item \textbf{Taylor keeps a single bulk momentum equation and adds molecular
    and eddy diffusion}, which produces the same qualitative diffusive
    separation at far lower cost. Two numbers are quoted in Taylor (2025) and
    they belong to different models: the best-fit model's elemental He/H falls
    from 8\,\% at the base of the thermosphere to about 2.5\,\% at high
    altitude, while the model Taylor sets against Xing's and Schulik \& Owen's
    (2025) multi-fluid results is Model~B, whose H/He goes from 92/8 at the
    base to about 96/4.
  \item \textbf{\AIOLOS\ is multi-fluid but escape-agnostic}: its friction
    solver would in principle give fractionation, but the code is not tuned for
    the transonic-wind problem.
  \item \textbf{\EXHALE\ keeps one bulk velocity and diffuses the composition
    on top of it}, which is Taylor's choice of architecture but not Taylor's
    variable. Where Taylor transports the trace ratio $n_{\rm He}/n_{\rm H}$,
    the binary operator
    (\texttt{binary\_element\_diffusion.f90}) transports the helium
    \emph{mass fraction} $X = \rho_{\rm He}/\rho$ of a two-component mixture,
    which is bounded at both ends of the composition axis and therefore does
    not break down as helium approaches unity. Hydrogen's carriers are
    H\,\textsc{i}, H\,\textsc{ii}, H$_2$, H$_2^+$ and H$_3^+$, helium's are
    He\,\textsc{i}, He\,\textsc{ii} and He\,\textsc{iii}, and the binary
    coefficient is a Blanc's-law mixture over their stage fractions, with
    hard-sphere (Banks \& Kockarts 1973), polarization ion--neutral (Schunk \&
    Nagy) and Coulomb ion--ion (Paquette et al.\ 1986) pair coefficients chosen
    by charge. That last point is the physical content: an ion is held to the
    protons by Coulomb friction instead of settling at a neutral rate.
\end{itemize}

\subsection{Radiation transport}
\AIOLOS\ is the only one with a full \textbf{FLD} thermal-radiation solver,
which is necessary for optically-thick interiors and dust and overkill for the
escaping thermosphere. Taylor is the only one doing \textbf{\Lya\ Monte Carlo}
transfer iterated with the hydro solution. Xing and \EXHALE\ both use radial
XUV attenuation: \EXHALE\ on an energy grid built per run from the
ionization thresholds, driven by either a power-law XUV spectrum or a
tabulated SED. On top of that \EXHALE\ closes \Lya\ trapping either with a
the escape probability of the static plane-parallel damping-wing slab
(Neufeld 1990 eq.~3.27, Harrington 1973 eq.~40) evaluated in line every step, or
with an externally computed $J_{\rm Ly\alpha}(r)$ profile read from file
(\texttt{radiation/lya\_rt.f90}; keys \texttt{Jlya escape-prob} and
\texttt{Jlya RT file}). The second of those is how a real Monte Carlo field
enters: the LaRT scattering-rate field is produced offline
(\texttt{examples/exhale\_to\_lart.py}) and read back in.

\subsection{Ionization and chemistry solver}
\label{sec:chemsolver}
Taylor uses the KPP kinetic preprocessor over a large stiff ODE network;
\AIOLOS\ uses the C2Ray two-species scheme of SB23, or the reaction network
its code has since grown; Xing folds
semi-implicit source terms into the PLUTO RK stages. \EXHALE\ solves a single
coupled \emph{nonlinear algebraic} system per cell, which is what a steady-state
code needs, and everything lives in it together: H, He, the He($2^3$S)
metastable, the ten metal elements, and the four molecular species. Adding
physics means adding rows to that system rather than a parallel decoupled
solver. The molecular rows are normalized by a turnover-rate scale before they
reach the solver, because the He and H blocks otherwise differ in magnitude by
six orders at He/H\,$=10^3$ and \texttt{hybrd1} does not scale rows itself.

\paragraph{Frelikh integrates the network in time; \EXHALE\ solves it for
equilibrium.} This is the sharpest methodological split in this memo, and it
is worth stating precisely because the two codes model the same layer. Frelikh
splits the fluid equations into an advective part and source terms, and solves
the chemistry in a substep of its own: of Equation~(6) they write that the
chemical source term is one ``which we solve for in a separate substep of our
routine.'' The network is then a stiff initial-value problem, and their
Section~2.9 says so:
\begin{quote}\small
``As the rates of the reactions vary by several orders of magnitude, the
reaction network cannot be solved via a straightforward explicit method---the
set of equations is said to be `stiff.' We implement the semi-implicit method
of solving stiff systems of equations from W.\ H.\ Press et al.\ (2007), the
details of which are given in Appendix B.''
\end{quote}
Their Appendix~B gives the whole of it: semi-implicit Euler as in Numerical
Recipes Section 17.5.2, the specific stepper named (\texttt{StepperSie}), the
linear system solved by LU decomposition with forward and back substitution,
$\Delta t$ crossed by a sequence of substeps $dt_i$ with Bulirsch--Stoer
polynomial extrapolation to test each one, and the tolerances quoted
(\texttt{rtol}~$=10^{-7}$, \texttt{atol}~$=10^{-4}\,$\texttt{rtol},
$h_1 = 10^{-6}$, \texttt{hmin}~$=0$). The Jacobian is analytic: ``The Jacobian
$\partial f/\partial y$ is calculated analytically for our problem, as are the
time derivatives of the concentrations on the left-hand side of
Equation~(B6).'' Cells are independent, so the network runs in parallel across
the grid.

Set against that, \EXHALE\ has the more careful \emph{formulation} and the
weaker \emph{linear algebra}. Its rows are steady-state production--loss
balances (\texttt{System\_HeH\_mol.f90}); there is no path in the tree that
integrates the network in time.
\texttt{nonlinear\_system\_solver/constrained\_chemical\_equilibrium.f90}
takes the unknowns to be $u_k = \ln n_k$, so no iterate can go negative,
following the CEA method of Gordon \& McBride (1994); it writes element
conservation as an explicit residual row per element rather than as a closure,
so that HeH$^+$ satisfies the H and the He budget simultaneously; and it takes
$n_e$ from charge neutrality rather than carrying it as an unknown, so
neutrality holds identically at every iterate. Each rung is then handed to
MINPACK \texttt{hybrd} with a \emph{forward-difference} Jacobian: the
opposite of Frelikh's choice.

The reason \EXHALE\ needs the extra apparatus is stated in that file's own
header: the molecular cell is bistable, with ``a molecular basin in the dense,
shielded base and an atomic basin above the \Htwo\ $\rightarrow$ H front,'' so
a single starting guess either lands in the right basin or does not. The fix
is a natural-parameter continuation in the radiation field, tracked from six
decades below the cell's own field (where the composition is known in
closed form) up to the true field, bisecting any rung that fails.

\textbf{The judgement worth recording is that Frelikh does not meet this
problem at all, and not because its solver is better.} A time-integrated
network has no uniqueness question in it: the basin is selected by the
previous step's state, not by a guess. Their only defense against a
composition leaving the physical domain is a hard floor, stated in their
Section~2.5 as ``We implement a number density floor of zero for all
species,'' and they record that the floor leaves a visible mark: their
Figure~4 caption attributes a kink in the \Htwo\ density profile to
``applying the density floor in a region where the profile drops off
steeply.'' \EXHALE's continuation is, in effect, a reconstruction inside an
equilibrium formulation of the path that time integration supplies for free.
Neither approach is wrong; the point is that the bistability is a cost of
choosing equilibrium, not a fact about the physics.

\subsection{The molecular reaction network}
\label{sec:network}

Frelikh's Table~1 and \EXHALE's \texttt{lower\_atmosphere/mol\_rates.f90} are
the only two molecular H--He networks among the five that can be compared
reaction by reaction, and they largely agree. Frelikh lists 22 reactions
compiled from Yelle (2004) and Garc\'ia-Mu\~noz (2007); \EXHALE\ carries 23
transcribed from Koskinen et al.\ (2022) Table~1, which was checked here
against the published table rather than against the transcription. All 22 of
Frelikh's arrows are one-directional; there is no reversible pair notation in
the table.

\textbf{Twelve are identical in value and in source}: H$^+$ and He$^+$
radiative recombination (Storey \& Hummer 1995), H$_2^+$ dissociative
recombination (Auerbach et al.\ 1977), H$_2^+ + $\Htwo\ (Theard \& Huntress
1974), H$_2^+ + $H (Karpas et al.\ 1979), H$^+ + $\Htwo\ (Yelle 2004),
three-body H$_2$ formation $8.0\times10^{-33}(300/T)^{0.6}$ (Ham et al.\
1970), He$^+ + $\Htwo\ $\rightarrow$ HeH$^+$ (Schauer et al.\ 1989),
HeH$^+ + $\Htwo\ (Bohme et al.\ 1980), HeH$^+ + $H (Karpas et al.\ 1979),
HeH$^+$ dissociative recombination (Yousif \& Mitchell 1989), and the
three-body \Hth\ association. The last of these is worth naming because it is
easy to assume it is missing: Frelikh's \textbf{k21},
H$^+ + $\Htwo$ + $M~$\rightarrow$~\Hth$ + $M at
$3.2\times10^{-29}$~cm$^6$\,s$^{-1}$ (Miller et al.\ 1968), \textbf{is present
in \EXHALE} as R13, \texttt{rk\_R13\_Hp\_H2\_M}, at the same value and from
the same source.

\textbf{Five differ}, and Table~\ref{tab:ratediff} gives the sizes. Four are
choices of source table. The fifth is not.

{\tabsetup
\begin{table}[htbp]\centering
\caption{Rate coefficients that differ between Frelikh's Table~1 and
\EXHALE's transcription of Koskinen et al.\ (2022) Table~1. Ratios evaluated
here.}\label{tab:ratediff}
\begin{tabular}{@{}L{0.20\textwidth} L{0.235\textwidth} L{0.235\textwidth}
                  L{0.20\textwidth}@{}}
\toprule
\textbf{Reaction} & \textbf{Frelikh} & \textbf{\EXHALE} & \textbf{Ratio} \\
\midrule
\Hth$+$e$\rightarrow$\Htwo$+$H &
$2.9{\times}10^{-8}(300/T)^{0.65}$ (Sundstr\"om 1994) &
$2.16{\times}10^{-8}(300/T)^{0.65}$ (Larsson 2008) & 1.34 \\
\Hth$+$e$\rightarrow$3H &
$8.6{\times}10^{-8}(300/T)^{0.65}$ (Datz 1995) &
$5.04{\times}10^{-8}(300/T)^{0.65}$ (Larsson 2008) & 1.71 \\
\emph{\Hth\ recombination, total} & $1.15{\times}10^{-7}(300/T)^{0.65}$ &
$7.20{\times}10^{-8}(300/T)^{0.65}$ & \textbf{1.60 at every $T$} \\
\Hth$+$H$\rightarrow$H$_2^+ +$\Htwo &
$2.0{\times}10^{-9}$, \textbf{no barrier} (Yelle 2004) &
$2.1{\times}10^{-9}e^{-20000/T}$ (Harada 2010) &
$4.6{\times}10^{8}$ (1000\,K); $4.6{\times}10^{6}$ (1300\,K);
$2.1{\times}10^{4}$ (2000\,K); 748 (3000\,K) \\
\Htwo$+$M$\rightarrow$H$+$H$+$M &
$1.5{\times}10^{-9}e^{-4.8\times10^{4}/T}$ &
$k_{\rm rec}/K_{\rm eq}$, detailed balance of R15
(\emph{since 2026-09-01; was $1.5{\times}10^{-9}e^{-48350/T}$, which both
codes cite from Baulch et al.\ 1992}) &
\EXHALE/Frelikh $= 0.079$ (1000\,K); 0.19 (1300\,K); 0.51 (2000\,K);
0.87 (3000\,K). Was 0.705 (1000\,K); 0.839 (2000\,K) \\
He$^+ +$\Htwo$\rightarrow$H$^+ +$H$+$He &
$8.8{\times}10^{-14}$ (Schauer 1989) &
$1.0{\times}10^{-9}e^{-5700/T}$ (Moses \& Bass 2000) &
38 (1000\,K); 142 (1300\,K); 657 (2000\,K) \\
\bottomrule
\end{tabular}
\end{table}}

\paragraph{\Hth$ + $H is a physics error in Frelikh, not a source choice.}
The reaction \Hth$ + $H~$\rightarrow$~H$_2^+ + $\Htwo\ is endothermic. From
the standard enthalpies of formation ($1488.3$~kJ\,mol$^{-1}$ for H$_2^+$
(which is $\mathrm{IE}(\Htwo) = 15.426$~eV), $1107$ for \Hth\ and $218.0$ for
H) the reaction enthalpy is
$\Delta H = 1488.3 - 1107 - 218.0 = 163.3$~kJ\,mol$^{-1} = 1.69$~eV, i.e.\ a
barrier of $19\,600$~K. The Harada et al.\ (2010) form that \EXHALE\ carries,
$e^{-20000/T}$, matches that endothermicity; a temperature-independent
$2.0\times10^{-9}$ does not, and there is no $v$-excited label on the entry as
there is on the neighbouring H$^+ + \Htwo(v{\ge}4)$ row. \textbf{The judgement
is that \EXHALE\ is right here and Frelikh is not}, at the four to eight
orders of magnitude the table gives for the molecular layer. The caveat is
that Yelle (2004) is not in \texttt{references/}, so whether the printed
constant faithfully transcribes its source could not be checked.

\textbf{Nothing in Frelikh's network is absent from \EXHALE.} Its k1, k3, k4,
k14 and k22 are photoreactions, and \EXHALE\ has H, He and \Htwo\
photoionization (\texttt{util\_ion\_eq.f90}) and Lyman--Werner
photodissociation (\texttt{lyman\_werner.f90}, opt-in) of its own. The one
unchecked item is whether \EXHALE\ implements the dissociative photoionization
channel \Htwo$ + h\nu \rightarrow$ H$ + $H$^+ + $e$^-$ separately.

\textbf{Six \EXHALE\ reactions have no counterpart in Frelikh}: collisional
ionization of H and of He (Voronov 1997), electron-impact dissociation of
\Htwo\ (Stibbe \& Tennyson 1999), \Htwo$ + $He$^+$ (Barlow 1984), and the two
He$\leftrightarrow$H charge-exchange channels. The last of these is the one
that matters. Frelikh's seven helium reactions (k14--k20) cover
photoionization, two He$^+ + $\Htwo\ channels, two HeH$^+$ channels,
He$^+$ recombination and HeH$^+$ dissociative recombination, and
\textbf{contain no charge exchange with hydrogen in either direction}. In
\EXHALE\ the live definitions are in
\texttt{radiation/charge\_exchange.f90}, Group~B, and the two are not each
other's reverse: B1, He$ + $H$^+ \rightarrow$ He$^+ + $H, is the
non-radiative collisional channel (Kimura et al.\ 1993), while B2,
He$^+ + $H~$\rightarrow$~He$ + $H$^+ + $photon, is radiative charge transfer
(Stancil et al.\ 1998). A photon-emitting process has no collisional reverse,
which is why the module forbids testing the pair against detailed balance.

\subsection{\texorpdfstring{\Hth}{H3+} infrared cooling}
\label{sec:h3p}

Both codes take the cooling function from Miller, Stallard, Tennyson \& Melin
(2013), and \EXHALE's 800--1800~K coefficients are the same seven numbers as
Frelikh's Equation~(13). Three things differ.

First, \textbf{Frelikh radiates into vacuum and \EXHALE\ can exchange with a
field}. Frelikh's volumetric rate is
$\Lambda_{\Hth} = 4\pi E(T)\,n_{\Hth}\times10^{7}$~erg\,cm$^{-3}$\,s$^{-1}$,
with no incident term, no stimulated emission and no optical-depth factor;
downward emission is treated as lost: ``We treat the radiation that is emitted
downward (i.e., toward the lower simulation boundary) as having left the
domain.'' Optical thinness is verified after the fact in their Section~3.4
against ExoCross cross sections. \EXHALE's
\texttt{lower\_atmosphere/h3p\_cooling.f90} carries the same emission term but
its \texttt{h3p\_net\_cooling\_rate} returns
$\Lambda_{\rm net} = n_{\Hth}\,4\pi\,s(T,n_{\Htwo})\,[\,E(T) - W_{\rm
dil}E(T_{\rm rad})\,]\times10^{7}$, which is exactly zero when $T = T_{\rm
rad}$ and $W_{\rm dil} = 1$. \textbf{But $W_{\rm dil}$ is reached only through
the key \texttt{Base IR field}, which defaults to off}, so \EXHALE's default
behavior is Frelikh's.

Second, \EXHALE\ carries the \textbf{non-LTE departure factor}
$s(T, n_{\Htwo})$ of Miller et al.'s Table~6 and all four published
temperature segments (30--300, 300--800, 800--1800, 1800--5000~K); Frelikh is
pure LTE and implements only the upper two, so its \Hth\ cooling is undefined
below 800~K and is set to zero above 5000~K.

Third, and recorded here because it will confuse anyone reading the paper
against the code: \textbf{Frelikh's Equations~(13) and~(14) are printed
without their logarithm}, as
$E(800 < T < 1800\,\mathrm{K}) = -62.7016 + 0.0526104\,T - \ldots$ carrying
units of W\,sr$^{-1}$\,molecule$^{-1}$. A negative polynomial cannot be a
power, and the missing operator is a \emph{natural} logarithm, not a common
one. Two independent checks were run here. The two published segments
disagree at their 1800~K junction by $-2.35\,\%$ read as $\ln E$ and by
$-5.34\,\%$ read as $\log_{10}E$, and \EXHALE's transcription of the same
Table~5 records the mismatch as ``2.4\,\% below'': the natural-log value.
And a two-level estimate for the $\nu_2$ band ($A \approx 100$~s$^{-1}$,
$E(\nu_2) = 3628$~K, $h\nu$ at 4~$\mu$m) gives
$1.05\times10^{-20}$~W\,sr$^{-1}$\,molecule$^{-1}$ at 1000~K, against
$3.34\times10^{-20}$ for the natural-log reading and $1.4\times10^{-45}$ for
the common-log reading. Their own figures confirm that the code behind the
paper is right and only the printed equations are wrong: their Figure~5 shows
$\Lambda_{\Hth}\approx2\times10^{-7}$~erg\,cm$^{-3}$\,s$^{-1}$ where Figure~4
shows $n_{\Hth}\approx$ a few$\times10^{3}$~cm$^{-3}$, which is the
natural-log ratio.

\subsection{Photoelectron secondary ionization}
\label{sec:secion}

Section~\ref{sec:common} records that \EXHALE\ no longer treats the
photoelectron heating fraction as a free parameter, which separates it from
Taylor and Xing. Frelikh does not treat it as free either, and it uses a
different and (for a molecular gas) a better prescription.

\EXHALE\ uses \textbf{Shull \& van Steenberg (1985)}, with the branching
between the H\,\textsc{i} and He\,\textsc{i} secondary channels renormalized
from the $n(\mathrm{He})/n(\mathrm{H}) = 0.1$ composition their Monte Carlo
sampled onto the cell's own neutrals (\texttt{svs85\_secondary\_branching},
\texttt{radiation/util\_ion\_eq.f90}). Frelikh uses \textbf{Dalgarno, Yan \&
Liu (1999)}: $W = W_0(1 + Cx^{\alpha})$ from their Tables 4 and 5, a secondary
ionization added to the network as an extra reaction, and (this is the
part \EXHALE\ has no equivalent of) the \Htwo\ channels put back
explicitly. Their Equations (25)--(35) return the 5.4~eV dissociation heat
following triplet excitation, the thermalized $\nu = 1$ (0.516~eV) and
$\nu = 2$ (1.032~eV) vibrational levels with three separate ionization-fraction
weightings, and the Lyman and Werner electronic states with their
$\sim$90\,\% fluorescent / $\sim$10\,\% dissociative branching.

Frelikh also states the range of validity and checks it: ``the approximations
presented in A.\ Dalgarno et al.\ (1999) are only valid for a gas with an
ionization fraction $<$0.1,'' which they justify by locating the relevant
$\tau_\nu = 1$ surfaces in the neutral layer, with a stated error of at most
15\,\%.

\textbf{The judgement is that for a molecular layer Dalgarno et al.\ (1999) is
the correct prescription and Shull \& van Steenberg (1985) is not.} SvS85 is a
Monte Carlo fit for a pure atomic H\,$+$\,He plasma; it carries no information
about where photoelectron energy goes in a gas that is mostly \Htwo. Checked
here, \texttt{svs85\_secondary\_branching} contains no \Htwo\ term at all:
the code carries an \Htwo\ photoabsorption column separately, but the
photoelectron energy partition has no \Htwo\ vibrational thermalization, no
dissociation heat return and no \Htwo\ secondary ionization. The two
prescriptions are complementary rather than competing: SvS85 remains valid to
high ionization fraction, where Dalgarno's fits are not, and Dalgarno covers
the molecular gas, where SvS85 has nothing to say. Neither covers the whole
domain.

The sign of the effect is also worth recording, since it is not the intuitive
one. Frelikh's Section~6 reports that including secondary ionization raised
the ionization fraction in the molecular layer, increased electron-impact
destruction of \Hth, reduced \Hth\ cooling and so \emph{increased} the escape
rate, ``from $6.4\times10^9$ to $9\times10^9$\,g\,s$^{-1}$\,sr$^{-1}$ for our
fiducial hot Jupiter run, which is counter to the intuition that secondary
ionization results in a decreased efficiency of escape.''

\subsection{Metals and cooling}
\EXHALE\ has the richest \emph{closed-form} metal treatment: ten elements
(C, O, N, Mg, Si, Ca, Na, K, S, Fe) in 27 ion stages inside the coupled
system, CHIANTI-fitted line cooling, ground-term fine-structure levels solved
in statistical equilibrium at the local $(n_e, n_{\rm HI})$ rather than in the
coronal limit, and line trapping in the eight fine-structure lines of
C\,\textsc{i}, C\,\textsc{ii}, N\,\textsc{ii} and O\,\textsc{i}. Taylor can
include solar-abundance metals but treats the structure as H/He-dominated.
Xing has \textbf{no metals} and \textbf{only \Lya\ cooling}. Frelikh has no
metals either, and only \Lya\ and \Hth\ cooling; it checked recombination
cooling in postprocessing against the Osterbrock \& Ferland rate and left it
off as negligible. Its Section~5.4 names the omission as the main limit on
extending the model to smaller planets, since \Hth\ is destroyed by H$_2$O and
CO and metals bring line coolants of their own. SB23 itself
carries only recombination and \Lya\ cooling; the \AIOLOS\ code adds the Black
(1981) rates with free-free and H$_3^+$ cooling, but neither is specialized to
the metal-line escape problem.

\subsection{Line diagnostics}
Taylor computes 10830 and \Ha\ in loop with the most complete non-LTE
microphysics. Xing post-processes 10830 only. \EXHALE\ builds the $2^3$S and
H($n{=}2$) populations in the run (the metastable is on by default), and
produces the transmission spectrum (10830, \Lya, \Ha, H$\beta$, the Mg\,II,
Ca\,II and Na\,I doublets) in \texttt{EXHALE\_transit.py}, with
impact-parameter Voigt integration, instrument and rotation convolution and an
optional triaxial Roche geometry. \AIOLOS\ does not target these lines.

\subsection{Lower boundary}
\label{sec:lowerbc}
Taylor couples to a photochemical lower and middle atmosphere at
$\sim$1~$\mu$bar, which is the most physically motivated boundary of the five.
Xing uses a fixed-density, fixed-temperature base. \AIOLOS\ leaves the boundary
configurable. \EXHALE\ now has the same range as Taylor: a fixed base, an
analytic chemical-equilibrium column, a scalar handoff, or a full profile,
item (D) below.

\paragraph{Frelikh imposes the same count as \EXHALE, and one thing more.}
Its base sits at $P_0 = 0.96$~dyn\,cm$^{-2}$, i.e.\ at ``the region of the
outflow above the 1~$\mu$bar level, near which molecules such as \Htwo\ are
present,'' with $T_0 = 1000$~K and $\Rp = 10^{10}$~cm for the fiducial hot
Jupiter (their Table~2). Their Section~2.8 states what is imposed and why:
\begin{quote}\small
``At the inner boundary, the flow is subsonic, so we specify the boundary
conditions for the number densities of each species and the pressure. We
extrapolate the velocity from the grid: $u_0 = u_1 - (r_1-r_0)(u_2-u_1)/
(r_2-r_1)$, a wave can travel from the domain to the inner boundary, as the
flow is subsonic, so the value of the velocity is part of the solution and
cannot be imposed as a boundary condition.''
\end{quote}
\EXHALE's \texttt{BC\_component\_constrho} pins $\rho$ and $p$ and copies the
ghost velocity from the first interior cell. \textbf{The count of imposed
conditions and the choice of which component floats are identical}, and
Frelikh gives the characteristic argument for that choice explicitly.
\textbf{The difference is that Frelikh also imposes the composition.} That is
the one item on \EXHALE's own outstanding list for the base that Frelikh has
already done; it appears as item~(M) in Section~\ref{sec:adopt}.

Two further points. Frelikh nowhere reports a supersonic base, nor a guard
against one. Searching the paper for that failure mode returns nothing: the
subsonic condition is a \emph{premise} of the Section~2.8 argument and the
possibility of its being violated is not discussed. The structural weakness
that follows (once the inflow is supersonic all three characteristics enter
the domain and an extrapolated velocity stops being a boundary condition)
applies to both codes; only \EXHALE\ has found it.

Second, Frelikh's answer to the base thermal problem is not a physical model
but a deliberate override, and the paper says so. Their Section~2.2 opens by
naming the symptom: ``We implement a bolometric heating and cooling term in
this region, as the EUV heating is insufficient to prevent the gas from
cooling to unphysically low temperatures'', and then applies, below the
$\tau = 3$ surface of the highest-energy bin,
$\Gamma_{\rm bol} - \Lambda_{\rm IR} = \kappa\sigma_{\rm SB}(T_0^4 - T^4)$
with the opacity set by hand:
\begin{quote}\small
``where we have set $\kappa_{\rm opt}$ and $\kappa_{\rm IR}$ to a single
inflated value $\kappa = 1$\,cm$^2$\,g$^{-1}$. In addition, in our
implementation, we set $T_0^4 = 4T_{\rm eq}^4$, such that Equation~(12) has
the form of a Newtonian cooling term. \ldots\ our choices effectively force
the bolometrically heated layer to be isothermal at the temperature given by
the input parameter $T_0$.''
\end{quote}
For the fiducial planet that region is $r < 1.005\,\Rp$. Heat conduction plays
no part in it: Frelikh's Section~3.5 finds conduction ``subdominant at all
fluxes in our grid of models,'' worth $\sim$2\,\% in $\Mdot$ at 0.3~au, and
the only stability remark about it in the paper runs the other way, their
Section~2.8 warns that ``conduction makes the choice of outer boundary
conditions more important, as the conduction term does not explicitly include
a flux limiter.''

The bearing on \EXHALE\ is that \EXHALE\ already owns a less blunt instrument
for the same problem. \texttt{Base IR field} gives the \Hth\ band a radiative
equilibrium fixed point at $T_{\rm rad} = 1140$~K, $W_{\rm dil} = 1/2$, i.e.\
975.4~K, instead of letting it run the layer down; it is off by default
(Section~\ref{sec:h3p}). The reservation is that the temperature trough
measured in the molecular gate cases reaches 865~K, \emph{below} that fixed
point, so \Hth\ alone does not account for it and the other channels in that
region have not been separated.

\subsection{The molecular layer, measured side by side}
\label{sec:mollayer}

Frelikh's fiducial hot Jupiter and \EXHALE's molecular regression cases are
close enough in irradiation to compare directly. Frelikh quotes
$F_{\rm EUV} = 2.6\times10^3$~erg\,cm$^{-2}$\,s$^{-1}$ at 0.05~au;
\texttt{EXHALE\_setup.out} for the gate cases reports
$\log_{10}F_{\rm XUV} = 3.193$, i.e.\
$1.56\times10^3$~erg\,cm$^{-2}$\,s$^{-1}$, a factor 1.7. The planets differ
more: 0.7~$M_{\rm J}$ against a 14.52~$M_\oplus$ hot Uranus
($\Rp = 0.49\,R_{\rm J}$, $T_{\rm eq} = 1140$~K, $a = 0.048$~au,
He/H~$= 0.0793$).

All \EXHALE\ numbers in Table~\ref{tab:mollayer} were measured here from
\texttt{backup/regression/*/output/} by parsing the \texttt{\# columns}
headers; no case was re-run. Frelikh's numbers are from its Table~2 and its
Figure~4, 5 and 9 captions, except the $\Lambda_{\Hth}$ plateau, which was
read off the Figure~5 axis and is therefore approximate.

{\tabsetup
\begin{table}[htbp]\centering
\caption{The molecular layer. \EXHALE\ columns measured here; Frelikh columns
as published.}\label{tab:mollayer}
\begin{tabular}{@{}L{0.20\textwidth} L{0.17\textwidth} L{0.15\textwidth}
                  L{0.17\textwidth} L{0.17\textwidth}@{}}
\toprule
& \textbf{Frelikh HJ} & \textbf{Frelikh s-Earth} &
\textbf{\EXHALE\ \texttt{mol\_metals}} &
\textbf{\EXHALE\ \texttt{mol\_base\_handoff}} \\
\midrule
Planet mass & 0.7\,$M_{\rm J}$ & 7.2\,$M_\oplus$ & 14.52\,$M_\oplus$ &
14.52\,$M_\oplus$ \\
Base pressure & 0.96\,$\mu$bar & 0.47\,$\mu$bar & 9.0\,$\mu$bar &
9.0\,$\mu$bar \\
Base $T$ & 1000\,K (forced) & 1000\,K (forced) & 1213\,K & 1213\,K \\
Base \Htwo\ fraction of H nuclei & $\approx1$ & $\approx1$ & 0.99989 &
0.9996 \\
\Htwo$\rightarrow$H crossover & 1.01\,$\Rp$ & no sharp transition &
1.158\,$\Rp$ (1166\,K) & 1.158\,$\Rp$ (1155\,K) \\
Molecular layer cut-off & 1.04\,$\Rp$, sharp & \textbf{none; \Htwo\ survives
throughout} & 1.19\,$\Rp$, sharp & 1.19\,$\Rp$, sharp \\
Molecular layer $T$ & 1000$\rightarrow$4000\,K & - &
1213$\rightarrow$1319\,K & 1213$\rightarrow$1305\,K \\
$T$ minimum / maximum & - / $\sim$9500\,K & - / 4200\,K &
\textbf{865\,K} at 1.190 / 2719\,K at 1.910 & 874\,K / 2647\,K \\
$n(\Hth)$ maximum & $\sim2{\times}10^{5}$\,cm$^{-3}$ & - &
46.3\,cm$^{-3}$ at 1.124\,$\Rp$ & $6.7{\times}10^{4}$\,cm$^{-3}$ at
1.031\,$\Rp$ \\
$N(\Hth)$ & $1.0{\times}10^{13}$\,cm$^{-2}$ &
$1.5{\times}10^{13}$\,cm$^{-2}$ & $1.73{\times}10^{10}$\,cm$^{-2}$ &
$1.79{\times}10^{13}$\,cm$^{-2}$ \\
$\Lambda_{\Hth}$ maximum & $\sim2{\times}10^{-7}$ (read off Fig.~5) & - &
$6.145{\times}10^{-10}$ at 1.113\,$\Rp$ & $8.898{\times}10^{-7}$ at
1.031\,$\Rp$ \\
\ldots as a fraction of local cooling & most of it & dominant & 6.8\,\% &
100\,\% \\
\ldots as a multiple of local heating & $\lesssim1$ & - & 0.003 &
\textbf{6.95} \\
Sonic point & 3.1\,$\Rp$ & 4.4\,$\Rp$ & 2.878\,$\Rp$ & 2.874\,$\Rp$ \\
Ionized fraction at sonic point & mostly ionized & mostly neutral & 0.41 &
0.41 \\
$\log_{10}\Mdot$ [g\,s$^{-1}$] & 11.05 if integrated over $4\pi$ & - &
10.582 & 10.575 \\
Peak Mach number below 1.35\,$\Rp$ & not reported & - & 0.017 & 0.014 \\
\bottomrule
\end{tabular}
\end{table}}

Three readings of that table matter.

\textbf{The three-layer structure reproduces, but on the wrong side of a
qualitative divide.} Frelikh's Section~5.4 sets its own super-Earth against
Koskinen et al.\ (2022) and finds them alike, and both unlike a hot Jupiter:
\begin{quote}\small
``Our super-Earth runs are qualitatively similar to the results of their hot
Uranus model, in which instead of the sharp transition between the molecular
hydrogen layer and the atomic layer, as seen in our hot Jupiter simulations
\ldots\ molecules are present throughout the outflow. Particularly interesting
is the presence of \Hth\ throughout the outflow, instead of being confined to
a narrow region at the base.''
\end{quote}
\EXHALE's hot Uranus gate does the opposite. Measured here, the \Htwo\
fraction of H nuclei falls from 0.48 at 1.16\,$\Rp$ to
$2.7\times10^{-5}$ at 1.30\,$\Rp$, and $n(\Hth)$ from 35 to
$1.2\times10^{-3}$\,cm$^{-3}$ over the same interval: four decades each.
The temperature is also low: 865--2719~K, against the 4000--5000~K Frelikh
quotes as the Koskinen hot Uranus prediction and the 1000--3800~K of its own
super-Earth. \textbf{\EXHALE's hot Uranus is behaving like Frelikh's hot
Jupiter, not like the Koskinen hot Uranus the gate case is meant to
reproduce.}

\textbf{The two \EXHALE\ gate cases disagree with each other by three orders
of magnitude} in $N(\Hth)$, $1.73\times10^{10}$ against
$1.79\times10^{13}$\,cm$^{-2}$, and only the second is in Frelikh's range.
The difference is the \texttt{base.inp} handoff, whose $q_{\rm H_2,base}$ sets
the base \Hth. In that case $\Lambda_{\Hth}$ is 6.95 times the local heating
at 1.031\,$\Rp$ and supplies 100\,\% of the cooling there, which is a net
cooling excess visible in the numbers.

\textbf{The base flow is not supersonic, and that is not the same as saying it
is healthy.} The peak Mach number below 1.35\,$\Rp$ is
$1.4$--$1.7\times10^{-2}$ in every gate case, and the He/H of these runs
(0.0793) is far outside the composition window where a supersonic base has
been seen. What the same output does show is described in
Section~\ref{sec:baseflaws}.

\section{The \Htwo\ thermal-dissociation pair: a defect diagnosed here, and
since repaired on one side}
\label{sec:detbal}

\textbf{Repaired on the \EXHALE\ side, 2026-09-01; see
Section~\ref{sec:detbalfix}.} The diagnosis below is left as it was written,
because it is what identified the defect and because it still holds for
Frelikh. The \EXHALE\ column of Table~\ref{tab:detbal} describes the code as it
was before that date.

Frelikh's k12 and k13 are the forward and reverse of the same process,
\Htwo$ + $M~$\rightleftharpoons$~H$ + $H$ + $M, taken from two different
measurements: $1.5\times10^{-9}e^{-4.8\times10^{4}/T}$ from Baulch et al.\
(1992) for the dissociation and
$8.0\times10^{-33}(300/T)^{0.6}$ from Ham et al.\ (1970) for the three-body
recombination. \EXHALE's R12 and R15 were the same pair from the same two
sources, differing only in that Koskinen et al.\ (2022) print the barrier as
$48\,350$~K where Frelikh rounds it to $4.8\times10^{4}$~K.

Nothing in Frelikh discusses whether the pair is thermodynamically consistent.
The paper was searched here for \emph{detailed balance}, \emph{thermodynamic},
\emph{equilibrium constant}, \emph{reverse reaction}, \emph{reversible},
\emph{extrapolat} and \emph{range of validity}, across the body, the
appendices, the footnotes and the table notes. The only hits are the linear
velocity extrapolation of Section~2.8 and the Bulirsch--Stoer extrapolation of
Appendix~B, both unrelated. \textbf{There is no statement about the
consistency of the pair, and none about the temperature range over which
either fit is valid.} No such record was found on the \EXHALE\ side either;
the caveats in \texttt{mol\_rates.f90} concern third-body efficiency and the
helium-rich limit, not forward--reverse consistency.

The ratio $k_{\rm diss}/k_{\rm rec}$ is the equilibrium constant
$n_{\rm H}^2/n_{\Htwo}$, with the third body cancelling, so it can be checked
against thermodynamics directly. Table~\ref{tab:detbal} does that.\footnote{%
$K = (2\pi\mu k_{\rm B}T/h^2)^{3/2}\,g_{\rm H}^2/z_{\rm int}(\Htwo)\;
e^{-D_0/k_{\rm B}T}$ with $\mu = m_{\rm H}/2$, $g_{\rm H} = 2$ and
$D_0 = 36\,118.11$~cm$^{-1}$. The internal partition function was summed
explicitly over $v = 0$--15 and $J = 0$--79 with the Huber \& Herzberg
constants ($\omega_e = 4401.21$, $\omega_ex_e = 121.33$,
$\omega_ey_e = 0.812$, $B_e = 60.853$, $\alpha_e = 3.062$~cm$^{-1}$),
discarding levels above $D_0$, with ortho/para nuclear-spin weights 3 and 1
divided out by $(2I+1)^2 = 4$ so that the nuclear spin cancels consistently
against the H atoms. A rigid-rotor harmonic-oscillator evaluation of the same
expression agrees to 2.8\,\% at 1000~K and 1.7\,\% at 2000~K.}

{\tabsetup
\begin{table}[htbp]\centering
\caption{The thermal-dissociation pair against thermodynamic equilibrium.
Ratios above unity mean the pair drives the gas more atomic than equilibrium
allows. Computed here. $^{\dagger}$\,\EXHALE\ as it was before 2026-09-01;
Section~\ref{sec:detbalfix}.}\label{tab:detbal}
\begin{tabular}{@{}r r r r r r@{}}
\toprule
$T$ [K] & $K_{\rm thermo}$ [cm$^{-3}$] & Frelikh $k_{12}/k_{13}$ & ratio &
\EXHALE\ R12/R15\rlap{$^{\dagger}$} & ratio \\
\midrule
 500 & $7.07\times10^{-22}$ & $5.17\times10^{-19}$ & \textbf{732} &
       $2.57\times10^{-19}$ & \textbf{363} \\
 800 & $7.71\times10^{-5}$ & $2.96\times10^{-3}$ & \textbf{38.4} &
       $1.91\times10^{-3}$ & \textbf{24.8} \\
1000 & $3.80\times10^{1}$ & $5.50\times10^{2}$ & \textbf{14.5} &
       $3.88\times10^{2}$ & \textbf{10.2} \\
1500 & $1.54\times10^{9}$ & $6.24\times10^{9}$ & 4.05 &
       $4.94\times10^{9}$ & 3.21 \\
2000 & $9.96\times10^{12}$ & $2.21\times10^{13}$ & 2.22 &
       $1.86\times10^{13}$ & 1.86 \\
3000 & $6.38\times10^{16}$ & $8.40\times10^{16}$ & 1.32 &
       $7.48\times10^{16}$ & 1.17 \\
4000 & $4.97\times10^{18}$ & $5.45\times10^{18}$ & 1.10 &
       $4.97\times10^{18}$ & 1.00 \\
5000 & $6.62\times10^{19}$ & $6.87\times10^{19}$ & 1.04 &
       $6.40\times10^{19}$ & 0.97 \\
\bottomrule
\end{tabular}
\end{table}}

\textbf{The judgement is that the pair is thermodynamically consistent only
above about 3000--4000~K, and that this is a defect of both codes equally.}
Below that it drives the composition toward atomic hydrogen faster than
equilibrium permits, by a factor 14.5 in Frelikh and 10.2 in \EXHALE\ at
1000~K, and by 38 and 25 at 800~K. \EXHALE\ is the less wrong of the two only
because the higher barrier it transcribes partly compensates; the error is the
same error. Both codes place their molecular layer in exactly the 1000--1500~K
band where the inconsistency is largest: Frelikh forces its base to
$T_0 = 1000$~K, and the \EXHALE\ gate cases measure 1213--1319~K through the
layer.

Two qualifications. Neither code sets the \Htwo/H balance from this pair
alone: photodissociation, the Lyman--Werner channel and the ion chemistry all
act, so the composition error is smaller than the tabulated factors and is
confined to wherever thermal dissociation and three-body recombination
dominate. And this is a statement about the pair, not about either published
measurement in its own range of validity; the Baulch fit belongs to
shock-tube temperatures far above the molecular base.

\subsection{The repair, 2026-09-01}
\label{sec:detbalfix}

\EXHALE\ no longer transcribes the dissociation fit.
\texttt{rk\_R12\_H2\_thdis} in
\texttt{src/modules/\allowbreak lower\_atmosphere/\allowbreak mol\_rates.f90}
now returns \texttt{k3b\_H\_H\_to\_H2(T)\allowbreak\ /\ \allowbreak
keq\_H\_H\_to\_H2(T)}: the three-body
recombination coefficient divided by the equilibrium constant the module
builds from the \Htwo\ spectroscopic constants. R15 moved at the same time
from the Ham et al.\ (1970) room-temperature measurement to the Cohen \&
Westberg (1983) recommended $k_1(\Htwo) = 2.8\times10^{-31}T^{-0.6}$, whose
stated 50--5000~K range covers the molecular layer where Ham's 77--300~K does
not.

The consequence for Table~\ref{tab:detbal} is that the \EXHALE\ ratio column
becomes \textbf{1 by construction at every temperature}: R12/R15 is now
$1/K_{\rm eq}$ exactly, and against the independently computed
$K_{\rm thermo}$ of the footnote it is 0.998--1.000 over 500--5000~K, the
residual being the two evaluations' different treatment of the \Htwo\ internal
partition function. The 363, 24.8, 10.2 and 3.21 of the 500--1500~K rows are
gone.

Two checks that the construction is right rather than merely self-consistent.
Against the Baulch fit it replaces, the new $k_{\rm diss}$ is 1.05 at 8000~K,
1.18 at 5000~K, 0.97 at 3000~K and 0.82 at 2500~K (agreement inside
Baulch's own stated 2500--8000~K range), and departs only below it (0.61 at
2000~K, 0.11 at 1000~K), which is where the fit is an extrapolation. Against
the primary shock-tube measurement rather than the evaluation, Breshears \&
Bird's own M~$=$~\Htwo\ coefficient $5.48\times10^{-9}e^{-52989/T}$, the ratio
is 1.11 at 3500~K (the bottom of their measured range), 0.99 at 4000~K, 0.82
at 5000~K and 0.51 at 8000~K.

\textbf{The judgement of this section therefore stands for Frelikh and no
longer for \EXHALE.} The pair is now thermodynamically exact in \EXHALE\ at
every temperature, and its molecular layer is no longer driven toward atomic
hydrogen faster than equilibrium permits. What the detailed-balance route
assumes instead is the vibrational state the coefficients refer to, which is
recorded at the code site and in
\texttt{docs/molecular\_hydrogen\_treatment.tex}.

\section{What the molecular gate cases show}
\label{sec:baseflaws}

The comparison in Section~\ref{sec:mollayer} required reading \EXHALE's
molecular regression output closely, and five things turned up in it that are
not about Frelikh at all. They are recorded here with the method, so that they
can be re-measured. Every one comes from
\texttt{backup/regression/\textless case\textgreater/} with no case re-run;
the profile quantities are from \texttt{output/Hydro\_ioniz.txt} columns
$r/\Rp$, $\rho$~[$m_{\rm H}$\,cm$^{-3}$] and $v$, and the warnings from
\texttt{run.log}. All five gate cases are relaxation snapshots pinned at
\texttt{maxsteps}~$=12\,000$.

\begin{enumerate}[nosep,leftmargin=*]
  \item \textbf{The mass flux is not flat at the base.} $\rho v r^2$ in
    \texttt{mol\_metals} runs $1.6563\times10^{16}$ at $r = 1.00019$ and
    $3.1420\times10^{16}$ at $r = 1.00039$: a factor \textbf{1.90 across
    one cell}, at the second interior cell. Over the molecular layer
    ($1.0 < r/\Rp < 1.19$) the spread $(\max-\min)/\mathrm{mean}$ is
    \textbf{1.865}. In \texttt{mol\_lyman\_werner} the first two rows are
    exactly zero, the third is \textbf{negative}, $-1.6507\times10^{16}$, and
    the fourth returns to $+1.3357\times10^{15}$: the flux changes sign inside
    the layer. A converged steady state requires $\rho v r^2$ to be constant
    in radius, and these are not.
  \item \textbf{Inflowing cells sit inside the molecular layer in one case.}
    Counting rows with $v < 0$: \texttt{mol\_metals},
    \texttt{mol\_base\_handoff} and \texttt{mol\_diffusion} each have 28, all
    at $r = 1.305$--$1.426\,\Rp$, i.e.\ \emph{outside} the molecular front at
    1.16--1.19. \texttt{mol\_lyman\_werner} has \textbf{144, at
    $r = 1.00019$--$1.12394\,\Rp$, inside the layer}, with the first cell's
    velocity exactly zero: the one-way valve latched shut.
  \item \textbf{The accepted chemical states that are not roots sit at the
    temperature minimum.} \texttt{mol\_metals/run.log} carries 24 lines of
    \texttt{(ioniz\_eq) NON-ROOT accepted (relaxation amnesty)} and
    \texttt{mol\_ir\_bands} 11; an example reads \texttt{cell 260 step 1 r
    1.197620 T[K] 1.197E+03 info 4 viol 1.210E-08 res 1.065E-03}. They cluster
    at $r \approx 1.19$--$1.20\,\Rp$, which is where the temperature trough
    bottoms out at 865~K (Table~\ref{tab:mollayer}).
    \textbf{The direction of causation was not determined.} Whether a runaway
    in the cooling breaks the chemical solve, or a state that is not a root
    drives the cooling, or a third cause produces both, is not decided by this
    measurement, and no inference is drawn here. What the coincidence does put
    in question is whether the thermal collapse and the non-uniqueness of the
    chemical solve are two problems or one.
  \item \textbf{Lyman--Werner self-shielding is being evaluated outside its
    fit.} \texttt{mol\_lyman\_werner/run.log}: ``\texttt{the H2 Lyman-Werner
    lines remove 1.01E+00 of the 912-1110 A band, i.e.\ all of it. The Draine
    \& Bertoldi (1996) self-shielding fit is a fit to a RATE and does not bound
    the photon budget; the band transmission is floored at zero and this column
    is outside the fit.}'' The band is 101\,\% removed and the transmission
    floors at zero. This is the same case as item~2; no connection between the
    two is asserted.
  \item \textbf{\texttt{mol\_ir\_bands} exercises one of the three channels its
    name implies.} With \texttt{Base IR field: True} and
    \texttt{Molecular IR bands: True}, the H$_2$O and CO infrared columns of
    \texttt{output/Cooling\_breakdown.txt} (col.\ 13 and 14) are
    \textbf{exactly zero at every radius}, because \texttt{Oxygen chemistry} is
    off in that configuration and the carriers do not exist. The only molecular
    infrared channel actually switched on is the \Htwo\ band, at
    $8.98\times10^{-9}$ at the \Hth\ peak cell and
    $2.93\times10^{-7}$~erg\,cm$^{-3}$\,s$^{-1}$ at its own maximum. As a
    regression gate the case does not protect the H$_2$O and CO paths.
\end{enumerate}

A note on quoting $\Mdot$ from these cases, since it is easy to get wrong.
\texttt{src/EXHALE\_main.f90} evaluates it at $j = N - 20$, i.e.\ at
$r = 3.9231\,\Rp$, and halves it under
\texttt{2D approximate method: Rate/2 + Mdot/2}. Recomputing
$4\pi\rho m_{\rm H} v r^2$ at that index and halving reproduces the
$\log_{10}\Mdot = 10.58$ that \texttt{EXHALE\_setup.out} prints. Evaluating at
the outer edge instead, or omitting the factor of two, gives 10.94 and is not
the quantity the code reports. In any case the mass-flux spread above means
none of these numbers should be quoted as a converged $\Mdot$.

For completeness: the temperature floor these cases could have hit is
$0.01\,T_0 = 11.4$~K, because \texttt{T\_trial} in
\texttt{energy\_semi\_implicit.f90} is dimensionless. No case activated it
(\texttt{grep -c floor} over \texttt{run.log} returns zero in all five), and
the measured minima (865, 874, 946 and 970~K) differ case by case rather
than resting on a common constant.

\section{Adoption assessment}
\label{sec:adopt}

Each candidate is given one of three verdicts. \textbf{Adopted} names the
module that implements it. \textbf{Adopted then replaced} records that the
element entered the code and was later superseded, and by what. \textbf{Not
adopted} states whether the original reason still holds.

\subsection{Adopted}

\paragraph{(A) He/H diffusive separation: adopted, in a stronger form than
recommended.} The recommendation was Taylor's route: a molecular-plus-eddy
diffusion term on top of a single bulk momentum solve. That is what
\texttt{functions/binary\_element\_diffusion.f90} does, with the two departures
described in Section~\ref{sec:fluid}, the transported variable is the helium mass
fraction of a binary mixture rather than a trace ratio, and the binary
coefficient is resolved by ionization stage rather than taken at a neutral
rate. The eddy term comes from the scalar key \texttt{He\_Kzz} or, when a
lower-atmosphere profile is in use, from that profile's $\Kzz(p)$ interpolated
onto the grid. One implicit backward-Euler step is taken per hydro step and
converged by Newton on a tridiagonal Jacobian, with a P\'eclet-switched
donor-cell treatment of the $X(1-X)$ drift product so the flux vanishes where
either component runs out. Metals ride with hydrogen at fixed metal/H unless
\texttt{He\_metal\_diffusion} is set, in which case each metal additionally
diffuses as a trace species through the solved hydrogen background. The
operator is available both inside the marching loop and, through an outer
composition relaxation, on the direct steady route. Default off
(\texttt{He\_diffusion}).

\paragraph{(D) Whole-atmosphere lower boundary: adopted, and carried past
Taylor.} \EXHALE\ now has three routes below the base, in increasing order of
what they carry. An analytic isothermal column from the 1-bar level to the
$\mu$bar base using the Koskinen et al.\ (2022, Eqs.~11--13) chemical-equilibrium
H$_2$/H/He partition (\texttt{lower\_atmosphere/lower\_column.f90}). A scalar
handoff file \texttt{base.inp}, generated analytically or from a VULCAN run,
which carries the EOS boundary ($T$, $r$, $p$, $q_{\rm H_2}$), the elemental
reservoirs (He/H and a \texttt{<El>\_H\_base} key for each of the ten
elements, which is what \texttt{melem\_ab} is then built from) and the eddy
coefficient. And a full \emph{profile} named by
\texttt{Lower atmosphere profile:}
(\texttt{files\_IO/lower\_atmosphere\_profile.f90}), which carries
$p$, $r$, $T$, $n_{\rm tot}$, $\rho$, $\Kzz$, $q_{\rm H_2}$, $q_{\rm H}$,
$X_{\rm He}$ and a mixing-ratio column per element, interpolated in $\log p$.
A scalar physics key sitting beside a profile is refused rather than merged,
and the two files must agree on a \texttt{solution\_id}.

The part that goes beyond Taylor is the closure. Taylor's coupling is one-way:
the photochemistry sets the base and the wind reads it. \EXHALE\ closes the
loop on the elemental fluxes. \texttt{binary\_element\_diffusion.f90} measures
$F_{\rm H}$ and $F_{\rm He}$ face by face and reduces them to a median and a
spread over a radial window, and \texttt{src/utils/element\_flux\_closure.py}
iterates $\Phi_{\rm El} = F_{\rm El}(\Phi_{\rm El})$ to a fixed point with
under-relaxation, the chemistry taking $\Phi$ as its upper boundary condition
and the wind returning $F$. The elemental composition of the wind is then an
output of the coupled pair rather than an input to either.
\texttt{photochem\_to\_lower\_profile.py} and
\texttt{vulcan\_to\_lower\_profile.py} produce the profile file; a residual
smaller than the spread of the window it was measured on is reported as
unresolved rather than as converged.

\paragraph{(E) Molecular chemistry and H$_3^+$ cooling: adopted.}
H$_2$, H$_2^+$, H$_3^+$ and HeH$^+$ are unknowns of the same coupled system as
H, He and the metals (\texttt{System\_HeH\_mol.f90},
\texttt{System\_HeH\_mol\_metals.f90}), so the molecular and metal blocks share
one electron density. H$_3^+$ infrared cooling follows Miller et al.\ (2013),
in LTE with the non-LTE departure factor and with absorption of the ambient
infrared field so the rate returned is the net one
(\texttt{lower\_atmosphere/h3p\_cooling.f90}). H$_2$ photodissociation in the
Lyman--Werner bands uses Draine \& Bertoldi (1996) self-shielding
(\texttt{lower\_atmosphere/lyman\_werner.f90}, key
\texttt{Stellar LW flux}, zero by default). The three-body density $M$ of the
association reactions is the total gas-particle density from
\texttt{calc\_ntot}, not a mass density. Default off
(\texttt{Molecular chemistry}); turning it on requires helium and forces the
molecular particle count at the base.

\paragraph{(F) \Lya\ Monte Carlo: adopted as an option, not as the default.}
The recommendation was to keep the escape-probability closure and use Taylor's
Monte Carlo only as a validation benchmark, because a WASP-121\,b column with
$\tau_0 = 1.2\times10^8$ makes an acceleration-free Monte Carlo intractable
and forbids core skipping. That still stands for the in-loop field: the default
is the in-line static-slab closure. But \EXHALE\ also accepts a
tabulated $J_{\rm Ly\alpha}(r)$ from any external solver
(\texttt{Jlya RT file}), and that path is used with the LaRT Monte Carlo code,
including the \Lya\ emitted in situ within the wind. The Monte Carlo is
therefore available where it is affordable, without being in the iteration.

\paragraph{(I) Navier--Stokes viscosity and heat conduction: adopted (from
Koskinen, i.e.\ Taylor's own hydro heritage).} This was not on the original
list. \texttt{time\_step/viscous\_conduction.f90} adds the viscous momentum
force, its dissipation and thermal conduction (the Navier--Stokes level that
CETIMB solves and that an inviscid HLLC scheme lacks) integrated
Crank--Nicolson and entering the steady residual through the same operator.
Both are opt-in (\texttt{Viscosity}, \texttt{Conduction}). Note that the module
does \emph{not} implement Koskinen et al.\ (2022) Eqs.~(B5) and (B6) as
printed: for a radial flow with the Newtonian stress
$\tau_{ij} = \mu(\partial_i w_j + \partial_j w_i) - \frac{2}{3}\mu
\delta_{ij}\nabla\!\cdot\!\mathbf{w}$, the viscous force and the dissipation
are
\[
  F_\mu = \frac{4}{3}\frac{1}{r^2}\frac{\partial}{\partial r}
          \!\left(r^2 \mu \frac{\partial w}{\partial r}\right)
        - \frac{4}{3}\frac{\partial \mu}{\partial r}\frac{w}{r}
        - \frac{8}{3}\frac{\mu w}{r^2},
  \qquad
  q_\mu = \frac{4}{3}\mu\!\left(\frac{\partial w}{\partial r}
          - \frac{w}{r}\right)^{\!2},
\]
and the code carries these rather than the printed forms; the derivation is in
the module header.

\paragraph{(J) \AIOLOS's opacity-model dispatcher: adopted.} Also not on the
original list, which treated \AIOLOS\ as a cross-check only.
\texttt{radiation/opacity\_models.f90} is an \AIOLOS-style dispatcher over
photoionization opacity models (analytic, constant, constant with
Robinson \& Catling pressure broadening, and tabulated from a two-column
\texttt{.opa} file) selected at runtime by the presence of
\texttt{opacity.inp}. The pressure factor that was a deferred hook in the
version it came through is applied cell by cell here.

\subsection{Adopted, then replaced}

\paragraph{(B) Temperature-dependent Penning ionization: adopted from Taylor
(2025), then retired.} The original recommendation was the cheapest item on the
list, and it was taken: the constant $5\times10^{-10}$~cm$^3$\,s$^{-1}$ of
Roberge \& Dalgarno (1982) was replaced by the two-branch power law of Taylor
et al.\ (2025) Table 2, unconditionally. It has since been withdrawn. The fit
was transcribed correctly but is not usable: it \textbf{jumps by a factor
1.5724 at 4000~K} ($1.5849\times10^{-9} \to 2.4921\times10^{-9}$), and a rate
coefficient is a continuous function of temperature. It also disagrees with its
own paper's Figure~19, where the Maxwell--Boltzmann average of the Cohen \&
Lane (1971) and Morgner \& Niehaus (1979) cross sections rises to
$1.50\times10^{-9}$ near 1400~K and falls to $1.29\times10^{-9}$ at 4000~K
while the tabulated branch decreases monotonically, and above 4000~K it exceeds
every cross-section determination collected by Garc\'ia Mu\~noz (2025).

What is in the code now is the continuous closed form of Garc\'ia Mu\~noz
(2025), \emph{A\&A} 698, A199: the Maxwell--Boltzmann average of the Movre \&
Meyer (1997) total ionization cross sections for the atomic partner
(\texttt{ioniz\_HeI23S\_H}) and of Cohen \& Lane (1977) for the molecular one
(\texttt{ioniz\_HeI23S\_H2}), both in \texttt{radiation/Cool\_coeff.f90}.
Both return the \emph{total} rate that removes the metastable; the 0.9/0.1
split between the Penning channel, which leaves a lasting ion and a free
electron, and the associative channel, which makes HeH$^+$, is carried by
\texttt{f\_penning\_HeI23S} at the points where the products matter. Taylor's
own Figure~19 agrees with the adopted expression to 10--20\,\%, so the two
independent calculations are consistent and it was the published fit that was
the outlier.

\paragraph{(K) \AIOLOS's C/N/O line-cooling fits: adopted, then replaced by
CHIANTI.} The analytic C/N/O metal-cooling fits of \AIOLOS's
\texttt{photochem.cpp} (code, not paper) were ported into
\texttt{Cool\_coeff.f90} and were the default for a time. They are now the legacy branch (\texttt{cno\_cool 0}) and
the default is a set of closed-form fits to CHIANTI, accurate to 0.2--2.1\,\%
over $10^3$--$10^5$~K. The reason is physical, not stylistic: the \AIOLOS\
fits contain no nitrogen cooling at all, and their O\,\textsc{i}/O\,\textsc{ii}
rates deviate from CHIANTI by 40--70\,\%. The \AIOLOS\ branch is kept so old
results remain reproducible, and its remaining tuning constants are flagged in
place.

\subsection{Not adopted}

\paragraph{(C) The B-spline K-matrix He($2^3$S) cross section, not adopted;
the reason still holds, but the earlier statement of it was too generous.} It
was claimed that \EXHALE\ was ``already most of the way there.'' What
\texttt{functions/cross\_sec.f90} actually has is a smooth background: two
Verner et al.\ (1996) single-shell forms joined by a log-linear bridge, fitted
at the few-percent level to the Norcross (1971) / TOPbase background, with the
autoionization region represented by a \emph{resonance-averaged} bump. Taylor's
cross section resolves the individual EUV resonances, which \EXHALE's does not
and cannot. The case for adopting it is unchanged: low cost, a data swap;
moderate value, since whether the resonances matter depends on where the
star's coronal lines fall, but it should be stated as a missing capability
rather than as a finished job.

\paragraph{(G) \AIOLOS's FLD radiation and general reaction network, not
adopted; the reason still holds, with one clause of it withdrawn.} FLD matters
for optically-thick interiors and dust, not for an optically-thin thermosphere,
and a general implicit reaction matrix is the wrong shape for a steady-state
code whose ionization system is one algebraic solve. Both still stand. The
third clause of the old argument (that HLLC would be no improvement on
\EXHALE's ``approximate Riemann solver'') was simply wrong about the code:
\EXHALE\ has had HLLC since ATES, alongside LLF and Roe, and HLLC is what every
example configuration uses. There was nothing to adopt there.

\paragraph{(H) Xing's electron-pressure Riemann split and full multi-fluid,
not adopted; the reason is now stronger.} Converting \EXHALE\ to genuine
multi-fluid means velocities per species, collisional drag and the
electron-pressure split, which is a rewrite of the hydro core and hard to
reconcile with a single coupled algebraic ionization solve. The original
argument was that Taylor's diffusion approximation captures most of the same
effect far more cheaply. That argument is now backed by what was built: the
stage-resolved binary operator reproduces the friction physics (an ion held
to the protons by Coulomb collisions, a neutral settling at a hard-sphere rate)
that motivated the multi-fluid treatment in the first place. Xing remains a
validation reference for the fractionation itself.

\paragraph{(L) The Dalgarno et al.\ (1999) photoelectron partition in a
molecular gas, not adopted; this one is a real gap.} Unlike (C), (G) and
(H), the reason here does \emph{not} still hold. \EXHALE\ uses Shull \& van
Steenberg (1985) throughout, and its branching routine contains no \Htwo\ term
of any kind (Section~\ref{sec:secion}): no vibrational thermalization, no
dissociation heat returned, no \Htwo\ secondary ionization. That was harmless
while the code had no molecular layer. It is not harmless now that
\texttt{Molecular chemistry} exists, because the layer it omits the physics
from is exactly the layer where the highest-energy photons deposit. The
prescriptions are complementary (SvS85 stays valid where Dalgarno's fits do
not, above an ionization fraction of 0.1), so the shape of the fix is a
composition switch between the two rather than a replacement. Cost: moderate,
a set of closed-form fits from Dalgarno et al.\ (1999) Tables 4 and 5 plus
their Equations~(25)--(35). Value: it changes the energy budget of the
molecular layer, and Frelikh reports the effect running counter to intuition,
raising $\Mdot$ rather than lowering it.

\paragraph{(M) A composition boundary condition at the base, not adopted;
recommended.} Frelikh imposes the number density of every species at the inner
boundary along with the pressure (Section~\ref{sec:lowerbc}); \EXHALE\ pins
$\rho$ and $p$ and lets the interior chemistry decide what the gas at the base
is made of. The imposed-condition count is the same either way, so this is not
a well-posedness question: it is a question of whether an inflowing
molecular reservoir is allowed to arrive with a composition of its own.
\EXHALE's own outstanding list already carries this item; the point worth
recording is that a peer code does it and reports no difficulty from it.

\paragraph{(N) Time integration of the network in molecular cells, not
adopted; recorded as an option rather than a recommendation.} The bistability
that forces the radiation-field continuation in
\texttt{constrained\_chemical\_\allowbreak equilibrium.f90} is a cost of
solving for
equilibrium, and a short time integration of the same network does not have
it (Section~\ref{sec:chemsolver}). Frelikh's stepper is an off-the-shelf
semi-implicit Euler with an analytic Jacobian, which is not expensive. Against
that: \EXHALE\ is a steady-state code, its molecular rows share one electron
density with the metal and He($2^3$S) rows, and splitting the molecular block
out into a time-integrated substep would break the single-system architecture
that the rest of this memo records as the right one. The honest statement is
that the continuation and the substep solve the same difficulty and neither
has been measured against the other here.

\paragraph{(O) Post-hoc verification of the radiative-transfer approximations,
not adopted; cheap, and worth doing.} Frelikh does two checks that
\EXHALE\ asserts rather than measures. Their Appendix~D runs a Monte Carlo of
1000 \Lya\ photons through the converged profile and reports 78\,\% escaping
to space, 11\,\% lost downward into the bolometric region and 11\,\%
thermalized, which is what licenses the optically thin \Lya\ cooling rate;
repeating it from deeper in the molecular layer gives 36/27/37 and they say so.
Their Section~3.4 checks \Hth\ optical thinness against ExoCross cross
sections at the computed column. \EXHALE\ has the ingredients for both (the
LaRT coupling for the first, the Miller et al.\ data for the second), and
uses neither as a check on its own closures.

\section{What was not verified}
\label{sec:unverified}

The Frelikh comparison rests in places on sources that were not available, and
on \EXHALE\ paths that were not followed to the end. Six items are open, and
are listed so that nothing above is read as more settled than it is.

\begin{enumerate}[nosep,leftmargin=*]
  \item \textbf{Yelle (2004) is not in \texttt{references/}}, so whether
    Frelikh's barrier-free \Hth$ + $H constant faithfully transcribes its
    stated source could not be checked. The thermochemical argument in
    Section~\ref{sec:network} stands on its own, but it convicts the printed
    rate, not necessarily the transcription.
  \item \textbf{Larsson et al.\ (2008) and Baulch et al.\ (1992) are likewise
    not held}, so the \Hth\ recombination factor of 1.60 and the $48\,350$
    versus $4.8\times10^{4}$~K barrier were compared between the two secondary
    sources, not against the primary measurements.
  \item \textbf{Whether \EXHALE\ implements the dissociative photoionization
    channel} \Htwo$ + h\nu \rightarrow$ H$ + $H$^+ + $e$^-$ was not
    established. \texttt{PIR\_H2} exists in \texttt{util\_ion\_eq.f90}; the
    branching among the \Htwo\ photo-channels was not traced.
  \item \textbf{The temperature clamp reported elsewhere was not reproduced.}
    The five molecular gate cases do not activate the 11.4~K floor, and their
    minima differ case by case (Section~\ref{sec:baseflaws}). Whatever run
    exhibited a clamped cell, it is not one of these.
  \item \textbf{The direction of causation between the temperature trough and
    the non-root chemical states} at $r \approx 1.19\,\Rp$ is undetermined
    (Section~\ref{sec:baseflaws}, item 3).
  \item \textbf{Frelikh contradicts itself on the number of evolved species},
    and the memo does not adjudicate. Its Section~2.9 says ``the seven
    atomic/molecular species ($+$electrons)''; its Appendix~A says ``the eight
    component species''; its Table~1 contains eight neutral and ionic species
    (H, H$^+$, \Htwo, H$_2^+$, \Hth, He, He$^+$, HeH$^+$).
\end{enumerate}

One further observation is recorded as an observation only. In Frelikh's own
molecular layer, reading $n(\mathrm{H}) \approx 10^{11}$ and
$n_e \approx 3\times10^{8}$~cm$^{-3}$ off its Figure~4, the barrier-free
\Hth$ + $H channel gives $\approx200$~s$^{-1}$ against $\approx12$~s$^{-1}$
for dissociative recombination, which would make it the dominant \Hth\ sink.
That sits awkwardly against their Section~2.3, which says of the
recombination channels that ``we verify post facto that they do indeed
dominate in our results.'' The numbers behind this were read off a published
figure, so it is not asserted as a finding.

\section{Summary}

\begin{itemize}[nosep]
  \item \textbf{\AIOLOS} is a general-purpose, multi-fluid RHD code with the
    most sophisticated \emph{numerics} and the least specialization to the
    He/\Ha\ escape-diagnostic problem. It has contributed less than Taylor to
    \EXHALE, but not nothing: the opacity-model dispatcher came from it, and
    its C/N/O cooling fits were \EXHALE's default until CHIANTI replaced them.
  \item \textbf{Taylor (2025/2026)} was and remains the closest match to
    \EXHALE's goals and its richest source of upgrades. Four of the six adopted
    items come from Taylor or from the Koskinen code beneath it: diffusive He/H
    separation, molecular chemistry with H$_3^+$ cooling, the whole-atmosphere
    lower boundary, and Navier--Stokes viscosity and conduction. The one
    Taylor item that did not survive is the Penning rate, and it failed on its
    own internal inconsistency rather than on anything about \EXHALE.
  \item \textbf{Xing (2023)} still provides the physically purest treatment of
    He/H mass fractionation, and is still not something to import. It is a
    validation reference.
  \item \textbf{Frelikh (2026)} is the closest peer to \EXHALE's molecular
    layer and, unlike the other three, it is a code \EXHALE\ can be checked
    \emph{against} rather than borrowed from: twelve of its 22 reactions are
    identical to \EXHALE's in value and source, and where the two differ
    \EXHALE\ is ahead on the network (it has the He$\leftrightarrow$H charge
    exchange Frelikh lacks, and the endothermic barrier on \Hth$ + $H that
    Frelikh's constant rate is missing), ahead on \Hth\ cooling (non-LTE
    departure factor, four temperature segments, an optional exchange with the
    base field), and behind on exactly two things: the photoelectron partition
    in a molecular gas, and a composition boundary condition at the base.
  \item \textbf{One finding belonged to no single code, and has since been
    repaired on one side.} The thermal-dissociation pair both codes took from
    Baulch (1992) and Ham (1970) is off thermodynamic equilibrium by a factor
    10--15 at 1000~K in the sense of over-dissociating
    (Section~\ref{sec:detbal}); \EXHALE\ now builds the dissociation from the
    recombination through the equilibrium constant and is exact
    (Section~\ref{sec:detbalfix}), and the finding stands for Frelikh only.
    The extrapolated-velocity inner boundary
    both codes use stops being a boundary condition if the inflow ever goes
    supersonic, a weakness only \EXHALE\ has so far found in itself, belongs
    to both.
  \item \textbf{The two gaps this memo was written to name are closed.}
    \EXHALE\ has radius-dependent He/H separation, and it has a lower boundary
    that a photochemical model sets rather than the user. What is left of the
    original list is one data swap (the resolved He($2^3$S) cross section) and
    two architectural items that were correctly declined. The Frelikh
    comparison has since opened four more, of which one (the missing \Htwo\
    channels in the photoelectron partition) is a genuine gap rather than a
    declined option.
\end{itemize}

\end{document}
